\section{Project Goal}

Lets now define the project's goal.

I want to compare politicians' speeches among three countries, namely Italy, France and the United States, between 1945 and 2025, and then cross-compare the resulting analyses with democracy level indicators taken from the V-Dem dataset\footnote{The V-Dem Dataset is a multidimensional and disaggregated dataset that reflects the complexity of the concept of democracy as a system of rule that goes beyond the simple presence of elections: \url{https://www.v-dem.net/data/the-v-dem-dataset/}.}.

Italy was a natural choice being the country I was born and raised in, therefore the one I am most knowledgeable of, but more than that it has always been a fascinating research case study, being often referred to as an "anomaly" or, by authors such as Marc Lazar, as a kind of \textit{laboratory} of a general change of European democracies. Since the 1990s, the state of Italian democracy has raised numerous questions, stimulating the curiosity of the international media, causing heated debates, providing material for seminars and conferences, and engendering academic and journalistic publications by Italians and foreigners alike. Lazar's opinion is that Italy is "at the same time an indicator of a more general democratic malaise and of the multiple attempts of invention of a new democratic order that are emerging all over the European Union". Similar phenomena affect, in their own ways, other European democracies, making it a meaningful case study \cite{lazar_testing_2013}.  

Italian, French and American histories are highly intertwined: throughout the Cold War, the United States encountered unexpected challenges from Italy and France, two countries with the strongest, and determinedly most anti-American, Communist Parties in Western Europe. Brogi \cite{brogi_confronting_2011} shows that the resistance to Americanization was a critical test for the French and Italian communists' own legitimacy and existence. Their anti-Americanism was mostly dogmatic and driven by the Soviet Union, but it was also, at crucial times, subtle and ambivalent, nurturing fascination with the American culture of dissent. The staunchly anticommunist United States, Brogi argues, found a successful balance to fighting the communist threat in France and Italy by employing diplomacy and fostering instances of mild dissent in both countries.

Relationship between these three countries was also one of collaboration: 

The three countries

for example the U.S. interjected in their mediterranean strategy \cite{brogi_competing_2006} and fought with the two countries communist parties, two of the strongest in Europe, throughout the second half of the $20^{th}$ century \cite{brogi_confronting_2011}. 



Moreover, the three countries share 


The first step is building a corpora of politicians speeches, one corpus for italian speeches, one for french speeches and one for american speeches. 